\section{Conclusion}

This paper revisited drift free repetitive motion control of redundant manipulators from a sampled control viewpoint. The original \methodone{}, \methodtwo{}, and \methodthree{} were first used to separate formulation preserving performance from objective reformulation performance. On this basis, \methodtwoDL{} was developed as the main method by tracking the sampled KKT residual of the classic \methodtwo{} formulation through projected DLCCZNN dynamics.

The main finding is specific and practically meaningful. If the classic equality constrained drift free formulation is retained, then \methodtwoDL{} is the most suitable main route among the solver transfers studied in this work. In the simulator based real time experiments, the tuned proposed method improves the original \methodtwo{} by about $33.1\%$ in mean position error and about $68.6\%$ in final drift norm. The tuning study further shows that final parameter selection cannot be made from offline replay alone; it must be guided by drift recovery under sampled execution. At the same time, the auxiliary studies show that algorithmic and structural decoupling can fail badly when the solver does not match the control structure, as illustrated by \methodthreelcc{}.

The paper also keeps an honest perspective on the secondary results. The Method~3 line, especially \methodthreepcg{}, can achieve even higher raw accuracy under a different discrete objective, but that line addresses a different design question from the main contribution here. For the present Control Engineering Practice target, the central story is therefore not that one solver dominates everything, but that sampled drift free control requires tight algorithmic and structural coupling, and that \methodtwoDL{} provides the clearest formulation preserving improvement.

Before submission, physical UR3e validation must be completed, repeated cycle hardware tests must be added, and the additional shape studies should be reorganized into a supplementary package. For a CEP version, the hardware experiment should appear as the culminating engineering validation rather than as a brief future work note.
